% AA-43 | Rozdział przygotowany ze starszego skryptu sources/pl/galois(2).tex.
% Dalsze ogniwa: AA-44 — twierdzenie Liouville'a, AA-45 — całki nieelementarne,
% AA-46 — podstawy różniczkowej teorii Galois.
\chapter{Ciała różniczkowe, derywacje i rozszerzenia różniczkowe}
\chaptermark{Ciała różniczkowe i derywacje}
\label{aa:ch:differential-fields}

Równanie algebraiczne można badać przez rozszerzenia ciał i ich
automorfizmy. Badanie równania $DF=f$, gdzie $D$ jest pochodną,
wymaga dodatkowo informacji o tym, jak różniczkowanie zachowuje się
przy dołączaniu nowych elementów. Zaczniemy od derywacji, ciał
stałych oraz rozszerzeń algebraicznych i przestępnych. Zbudowane
narzędzia posłużą w kolejnym rozdziale do sformułowania i dowodu
twierdzenia Liouville'a o całkowaniu w funkcjach elementarnych.

\section{Derywacje i ciała stałych}

\begin{definicja}[Ciało różniczkowe]
\label{aa:def:differential-field}
Niech $K$ będzie ciałem charakterystyki $0$.
\emph{Derywacją} na $K$ nazywamy odwzorowanie $D:K\to K$,
dla którego dla wszystkich $a,b\in K$ zachodzą
$D(a+b)=Da+Db$ oraz $D(ab)=aDb+bDa$.
Parę $(K,D)$ nazywamy \emph{ciałem różniczkowym}, a
\[
C_K=\{a\in K:Da=0\}
\]
jego \emph{ciałem stałych}.
\end{definicja}

\begin{stwierdzenie}[Reguły rachunku różniczkowego]
\label{aa:st:differential-field-rules}
Niech $(K,D)$ będzie ciałem różniczkowym. Wtedy
$D(1)=0$, $D(-a)=-Da$, $D(a^{-1})=-a^{-2}Da$
dla $a\ne0$ oraz
\[
D\!\left(\frac ab\right)=\frac{bDa-aDb}{b^2}
\quad(a,b\in K,\ b\ne0).
\]
Zbiór $C_K$ jest podciałem $K$, zawierającym $\mathbb Q$.
Dla każdej liczby całkowitej $n$ i $a\in K^\times$
zachodzi $D(a^n)=na^{n-1}Da$.
\end{stwierdzenie}
\begin{proof}
Z reguły iloczynu dla $1\cdot1=1$ wynika
$D1=2D1$, stąd $D1=0$. Addytywność daje
$D(-a)=-Da$. Jeśli $a\ne0$, to
$0=D(aa^{-1})=(Da)a^{-1}+aD(a^{-1})$,
co daje regułę odwrotności, a następnie ilorazu.
Jądro addytywnego odwzorowania $D$ jest zamknięte
na dodawanie, odejmowanie i mnożenie; reguła
odwrotności zapewnia też zamkniętość na odwrotności
elementów niezerowych. Każda liczba całkowita ma
pochodną zero, podobnie jak jej odwrotność w $K$,
więc $\mathbb Q\subseteq C_K$.
Regułę dla dodatnich $n$ otrzymujemy przez indukcję
z reguły iloczynu, a dla ujemnych z reguły
odwrotności. Przypadek $n=0$ został już wykazany.
\end{proof}

\begin{przyklad}[Funkcje wymierne]
\label{aa:prz:differential-rational-field}
Niech $C$ będzie ciałem charakterystyki $0$, a
$K=C(x)$ ciałem funkcji wymiernych jednej zmiennej.
Istnieje jedyna derywacja $D$ o własnościach
$D(c)=0$ dla $c\in C$ i $Dx=1$.
Dla $P,Q\in C[x]$, $Q\ne0$, reguła ilorazu przyjmuje postać
\[
D(P/Q)=\frac{P'Q-PQ'}{Q^2}.
\]
Pokażemy, że $C_K=C$. Jeśli niezerowa funkcja $P/Q$,
zapisana w postaci nieskracalnej, spełnia $D(P/Q)=0$,
to $P'Q=PQ'$. Względna pierwszość $P$ i $Q$
w $C[x]$ oznacza, że $P\mid P'$.
Gdyby $\deg P\ge1$, w charakterystyce $0$ mielibyśmy
$\deg P'=\deg P-1$, sprzeczność. Zatem $P$
jest stały; równość $P'Q=PQ'$ implikuje wtedy
$Q'=0$, czyli także $Q$ jest stały.
Funkcja zerowa oczywiście należy do $C$.
\end{przyklad}

\begin{przyklad}[Derywacja zerowa i rola charakterystyki]
\label{aa:prz:zero-derivation-and-positive-characteristic}
Derywacja $D=0$ na dowolnym ciele charakterystyki $0$
ma całe to ciało jako ciało stałych.
Nie można natomiast bez zmiany twierdzeń zastąpić
założenia o charakterystyce $0$ charakterystyką $p>0$:
w $\mathbb F_p(x)$, przy $Dx=1$, mamy
$D(x^p)=p x^{p-1}=0$, choć $x^p$ nie należy
do $\mathbb F_p$. Co więcej, dołączenie pierwiastka
$\alpha$ równania $\alpha^p=x$ nie pozwala przedłużyć
derywacji z $Dx=1$, bo
$1=Dx=D(\alpha^p)=p\alpha^{p-1}D\alpha=0$.
\end{przyklad}

\begin{definicja}[Rozszerzenie różniczkowe]
\label{aa:def:differential-extension}
Niech $(K,D)$ będzie ciałem różniczkowym.
\emph{Rozszerzeniem różniczkowym} $L/K$ nazywamy
rozszerzenie ciał wraz z derywacją, również
oznaczaną przez $D:L\to L$, której ograniczenie
do $K$ jest równe pierwotnej derywacji.
Mówimy, że jest ono \emph{bez nowych stałych},
jeśli $C_L=C_K$.
Homomorfizm $\sigma:L\to M$ rozszerzeń różniczkowych
nad $K$ nazywamy \emph{różniczkowym}, jeśli
$\sigma|_K=\mathrm{id}_K$ i $D\sigma(a)=\sigma(Da)$
dla każdego $a\in L$.
\end{definicja}

Warunek $C_L=C_K$ jest istotny, ponieważ z samej
własności $\sigma|_K=\mathrm{id}_K$ nie wynika,
że funkcja pierwotna ma jednoznaczną stałą całkowania
w $C_K$. Pokażemy niżej, jak zapewnić brak nowych
stałych w konstrukcjach logarytmicznych.

\section{Algebraiczne przedłużanie derywacji}

\begin{twierdzenie}[Jednoznaczne przedłużenie]
\label{aa:tw:algebraic-derivation-extension}
Niech $(K,D)$ będzie ciałem różniczkowym
charakterystyki $0$, a $L/K$ dowolnym rozszerzeniem
algebraicznym. Wówczas istnieje dokładnie jedna
derywacja $L\to L$ przedłużająca $D$.

Jeśli $\alpha\in L$ ma nad $K$ unormowany wielomian
minimalny $m_\alpha(T)\in K[T]$, a
$m_\alpha^D(T)$ powstaje przez zastosowanie $D$
do wszystkich jego współczynników, to
\[
D\alpha=-\frac{m_\alpha^D(\alpha)}{m_\alpha'(\alpha)}.
\]
Mianownik jest niezerowy.
\end{twierdzenie}
\begin{proof}
Każdy nierozkładalny wielomian nad ciałem charakterystyki
$0$ jest rozdzielczy (zob.
Wniosek~\ref{aa:wn:charzero-separable}). W szczególności
$m_\alpha'(\alpha)\ne0$.
Jeśli przedłużenie istnieje, zastosowanie derywacji
do równości $m_\alpha(\alpha)=0$ oraz reguły iloczynu
wymusza
$0=m_\alpha^D(\alpha)+m_\alpha'(\alpha)D\alpha$.
Otrzymujemy wskazany, jedyny możliwy wzór.

Dla ustalonej $\alpha$ zdefiniujmy $D\alpha$ tym wzorem.
Dla wielomianu $f\in K[T]$ połóżmy
$D(f(\alpha))=f^D(\alpha)+f'(\alpha)D\alpha$.
Musimy wykazać niezależność od przedstawienia.
Jeśli $f(\alpha)=g(\alpha)$, to
$f-g=m_\alpha h$ dla pewnego $h\in K[T]$.
Reguła Leibniza dla wielomianów daje
\[
(f-g)^D(\alpha)+(f-g)'(\alpha)D\alpha
=h(\alpha)\bigl(m_\alpha^D(\alpha)
                +m_\alpha'(\alpha)D\alpha\bigr)=0,
\]
ponieważ wszystkie wyrazy zawierające
$m_\alpha(\alpha)$ znikają. Operacja jest więc
dobrze określona na $K[\alpha]$.
Bezpośrednio z definicji i reguły iloczynu
dla wielomianów wynika addytywność i reguła
Leibniza. Ponieważ $\alpha$ jest algebraiczny,
$K[\alpha]=K(\alpha)$ jest ciałem; wzór na
odwrotność zapewnia derywację na całym $K(\alpha)$.

Jeżeli $L/K$ jest skończone, możemy dołączać
skończenie wiele generatorów algebraicznych
jeden po drugim. Na każdym piętrze przedłużenie
jest jednoznaczne, więc wynik nie zależy
od wybranej kolejności generatorów.
W rozszerzeniu dowolnym każdy element $L$
należy do skończonego podrozszerzenia.
Na sumie dwóch takich podrozszerzeń
przedłużenia z poprzedniego kroku są jednoznaczne;
zgadzają się więc na ich części wspólnej.
Definiując pochodną każdego elementu przez
jego dowolne skończone podrozszerzenie,
otrzymujemy spójną derywację na $L$.
Jednoznaczność dla $L$ wynika z tego samego wzoru.
\end{proof}

\begin{wniosek}[Brak nowych stałych przy rozszerzeniu algebraicznym]
\label{aa:cor:algebraic-constants}
Niech $(K,D)$ będzie ciałem różniczkowym
charakterystyki $0$, którego ciało stałych $C=C_K$
jest algebraicznie domknięte. Wtedy dowolne
rozszerzenie algebraiczne $L/K$, wyposażone
w jedyne przedłużenie derywacji, spełnia $C_L=C$.
\end{wniosek}
\begin{proof}
Oczywiście $C\subseteq C_L$. Niech $\alpha\in C_L$,
a $m_\alpha\in K[T]$ będzie unormowanym
wielomianem minimalnym. Z poprzedniego twierdzenia
i $D\alpha=0$ wynika $m_\alpha^D(\alpha)=0$.
Ponieważ współczynnik wiodący $m_\alpha$ wynosi $1$,
mamy $\deg m_\alpha^D<\deg m_\alpha$,
chyba że $m_\alpha^D=0$. Minimalność wielomianu
$m_\alpha$ wymusza $m_\alpha^D=0$.
Wszystkie współczynniki $m_\alpha$ należą zatem do $C$,
więc $\alpha$ jest algebraiczny nad $C$.
Algebraiczna domkniętość $C$ daje $\alpha\in C$.
\end{proof}

\begin{przyklad}[Pochodna funkcji algebraicznej]
\label{aa:prz:derivative-algebraic-sqrt}
Niech $K=\mathbb C(x)$ z derywacją $Dx=1$, a
$\alpha^2=x^2+1$ w rozszerzeniu algebraicznym.
Wielomian $T^2-(x^2+1)$ ma nad $K$ stopień $2$
i jest nierozkładalny: $x^2+1$ ma w $\mathbb C(x)$
proste zera $i$ i $-i$, więc nie jest kwadratem
funkcji wymiernej. Różniczkując relację, otrzymujemy
$2\alpha D\alpha=2x$, czyli
$D\alpha=x/\alpha$. Jest to także wynik wzoru
z poprzedniego twierdzenia.
\end{przyklad}

\begin{przyklad}[Dlaczego domkniętość algebraiczna stałych jest potrzebna?]
\label{aa:prz:algebraic-new-constants}
Niech $K=\mathbb Q(x)$, $Dx=1$, $Dq=0$
dla $q\in\mathbb Q$. Z przykładu
\ref{aa:prz:differential-rational-field} mamy
$C_K=\mathbb Q$, które nie jest algebraicznie
domknięte. W rozszerzeniu
$L=\mathbb Q(i)(x)$ jedyne przedłużenie
derywacji spełnia $D(i)=0$, bo
$i^2+1=0$ i $2iD(i)=0$. Dlatego
$C_L=\mathbb Q(i)\ne\mathbb Q$.
\end{przyklad}

\section{Ślad i norma a derywacja}

\begin{twierdzenie}[Różniczkowanie śladu i normy]
\label{aa:tw:differential-trace-norm}
Niech $(K,D)$ będzie ciałem różniczkowym
charakterystyki $0$, a $L/K$ skończonym rozszerzeniem
algebraicznym z jedynym przedłużeniem $D$.
Dla każdego $v\in L$ i $u\in L^\times$ zachodzi
\[
D\!\left(\operatorname{Tr}_{L/K}(v)\right)
=\operatorname{Tr}_{L/K}(Dv),
\qquad
\frac{D\!\left(N_{L/K}(u)\right)}{N_{L/K}(u)}
=\operatorname{Tr}_{L/K}\!\left(\frac{Du}{u}\right).
\]
Nie zakładamy normalności $L/K$.
\end{twierdzenie}
\begin{proof}
Wybierzmy domknięcie algebraiczne $\overline K$
z jedynym przedłużeniem $D$ z poprzedniego
twierdzenia. Ponieważ $\operatorname{char}K=0$,
rozszerzenie $L/K$ jest rozdzielcze.
Niech $\sigma_1,\ldots,\sigma_d:L\to\overline K$
będą wszystkimi $K$-zanurzeniami, gdzie $d=[L:K]$.
Dla każdego $w\in L$ mamy
$D(\sigma_j(w))=\sigma_j(Dw)$: istotnie,
obie strony są wymuszone przez wielomian minimalny
elementu $w$ nad $K$ i wzór na pochodną jego
pierwiastka. Korzystając ze wzorów na ślad
i normę przez zanurzenia
(Twierdzenie~\ref{aa:tw:trace-norm-embeddings}),
otrzymujemy
\[
\begin{aligned}
D\operatorname{Tr}_{L/K}(v)
 &=D\left(\sum_{j=1}^d\sigma_j(v)\right)
  =\sum_{j=1}^d\sigma_j(Dv)
  =\operatorname{Tr}_{L/K}(Dv),\\
\frac{DN_{L/K}(u)}{N_{L/K}(u)}
 &=\frac{D\prod_{j=1}^d\sigma_j(u)}
         {\prod_{j=1}^d\sigma_j(u)}
  =\sum_{j=1}^d\frac{\sigma_j(Du)}{\sigma_j(u)}
  =\operatorname{Tr}_{L/K}\left(\frac{Du}{u}\right).
\end{aligned}
\]
Wszystkie czynniki normy są niezerowe,
gdyż zanurzenia ciał są różnowartościowe.
\end{proof}

\begin{przyklad}[Ślad i norma w rozszerzeniu kwadratowym]
\label{aa:prz:differential-quadratic-trace-norm}
W przykładzie \ref{aa:prz:derivative-algebraic-sqrt}
mamy $\alpha^2=x^2+1$ i
$\operatorname{Tr}_{L/K}(\alpha)=0$,
$N_{L/K}(\alpha)=-(x^2+1)$.
Pierwsza tożsamość daje
$\operatorname{Tr}_{L/K}(D\alpha)=D0=0$.
Druga ma po lewej stronie
$D(-(x^2+1))/(-(x^2+1))=2x/(x^2+1)$.
Po prawej, ponieważ $D\alpha/\alpha=x/(x^2+1)\in K$,
otrzymujemy
$\operatorname{Tr}_{L/K}(D\alpha/\alpha)=2x/(x^2+1)$.
W ten sposób obie tożsamości można sprawdzić bezpośrednio.
\end{przyklad}

\section{Przestępne dołączanie funkcji pierwotnych}

Niech odtąd $(K,D)$ ma charakterystykę $0$ i
$t$ będzie \emph{przestępny} nad $K$.
Podając $Dt=a\in K$, rozszerzamy $D$ jednoznacznie
na $K(t)$ przez formalną regułę iloczynu i ilorazu.
W następnym lemacie ustalimy dokładnie, kiedy
takie dołączenie zachowuje ciało stałych.

\begin{lemat}[Kryterium braku nowych stałych]
\label{aa:lem:primitive-extension-constants}
Niech $(K,D)$ będzie ciałem różniczkowym
charakterystyki $0$, $C=C_K$, a $t$ będzie
przestępny nad $K$ i spełniał $Dt=a\in K$.
Wtedy $C_{K(t)}=C$ wtedy i tylko wtedy,
gdy nie istnieje $b\in K$ spełniający $Db=a$.
\end{lemat}
\begin{proof}
Jeżeli $Db=a$, to $D(t-b)=0$, a
$t-b\notin K$ z powodu przestępności $t$.
Powstaje zatem nowa stała.

W drugim kierunku przypuśćmy, że
$R\in K(t)\setminus C$ i $DR=0$.
Funkcję wymierną $R$ można rozwinąć
jednoznacznie jako formalny szereg Laurenta
w nieskończoności
$R=\sum_{j\le m}r_jt^j$, gdzie $r_j\in K$
i $r_m\ne0$, a $m\in\mathbb Z$.
Derywacja jest na tym szeregu zgodna z
\[
D(r_jt^j)=(Dr_j)t^j+j\,a\,r_jt^{j-1}.
\]
Tożsamość zachodzi również dla ujemnych $j$,
a współczynnik każdej potęgi pochodzi ze
skończenie wielu wyrazów.

Z porównania współczynników przy $t^m$
w równości $DR=0$ otrzymujemy $Dr_m=0$,
czyli $r_m\in C^\times$.
Jeżeli $m=0$, odejmujemy stałą $r_0\in C$.
Otrzymany niezerowy szereg ma najwyższy
wykładnik ujemny (gdyby był zerowy,
$R$ od początku należałoby do $C$).
Możemy więc w każdym przypadku założyć $m\ne0$.
Współczynnik przy $t^{m-1}$ daje
$Dr_{m-1}+ma r_m=0$; jeżeli w rozwinięciu
nie ma wyrazu $t^{m-1}$, przyjmujemy
$r_{m-1}=0$. Wobec tego
$b=-r_{m-1}/(mr_m)\in K$ spełnia $Db=a$.
To kończy dowód przez kontrapozycję.
\end{proof}

\begin{przyklad}[Dołączenie logarytmu]
\label{aa:prz:logarithmic-extension}
Weźmy $K=\mathbb C(x)$, $Dx=1$ oraz
$a=1/x$. W $K$ nie ma elementu $b$
spełniającego $Db=1/x$: po rozkładzie
funkcji wymiernej na ułamki proste
pochodna każdego składnika ma przy
każdym skończonym biegunie resztę $0$,
natomiast $1/x$ ma w $0$ resztę $1$.
Dołączając przestępny element $\ell$ z
$D\ell=1/x$, otrzymujemy rozszerzenie
różniczkowe $K(\ell)/K$ bez nowych stałych.

Lokalnie analitycznie $\ell$ może być
zarodkiem wybranej gałęzi funkcji $\log x$.
Uzasadnijmy również jej przestępność.
Gdyby taki zarodek był algebraiczny nad $K$,
to dla $L=K(\ell)$ i $d=[L:K]$ poprzednie
twierdzenie dałoby
\[
D\!\left(\frac{\operatorname{Tr}_{L/K}(\ell)}d\right)
=\frac{\operatorname{Tr}_{L/K}(D\ell)}d
=\frac{\operatorname{Tr}_{L/K}(1/x)}d
=\frac1x.
\]
Po lewej byłaby pochodna funkcji wymiernej
z $\mathbb C(x)$, co przeczy sprawdzonej
właściwości jej reszt. Zatem wybrana
lokalna gałąź logarytmu rzeczywiście jest
przestępna nad $\mathbb C(x)$.
\end{przyklad}

\begin{przyklad}[Sztuczne dołączenie nowej stałej]
\label{aa:prz:new-constant-obstruction}
Niech $K=\mathbb C(x)$, $Dx=1$ i dołączmy
przestępny element $t$ z $Dt=1$.
Ponieważ $Dx=1$, element $t-x$
jest nową stałą w $K(t)$.
Nie wolno więc automatycznie nazywać
dowolnej przestępnej zmiennej spełniającej
równanie pierwszego rzędu rozszerzeniem
bez nowych stałych.
\end{przyklad}

\section{Rozszerzenia wykładnicze i elementarne}

\begin{lemat}[Stałe w rozszerzeniu wykładniczym]
\label{aa:lem:exponential-extension-constants}
Niech $(K,D)$ będzie ciałem różniczkowym
charakterystyki $0$, $C=C_K$, a $t$ będzie
przestępny nad $K$ i spełniał $Dt=ht$
dla pewnego $h\in K$. Wtedy $C_{K(t)}=C$
wtedy i tylko wtedy, gdy dla żadnej liczby
całkowitej $m\ne0$ i żadnego $u\in K^\times$
nie zachodzi $Du/u=-mh$.
\end{lemat}
\begin{proof}
Jeśli taka para $(m,u)$ istnieje,
to $D(ut^m)=0$, a $ut^m\notin K$
z powodu przestępności $t$. Jest więc
nową stałą.

Odwrotnie, przypuśćmy, że
$R\in K(t)\setminus C$ i $DR=0$.
Rozwijamy $R$ formalnie w nieskończoności
jako $\sum_{j\le m}r_jt^j$ z $r_m\ne0$.
Tym razem
$D(r_jt^j)=(Dr_j+jhr_j)t^j$.
Jeśli $m=0$, z równania współczynnika
$t^0$ otrzymujemy $Dr_0=0$.
Odejmujemy $r_0\in C$; pozostały niezerowy
szereg ma najwyższą potęgę $m<0$.
Zatem możemy przyjąć $m\ne0$.
Współczynnik przy $t^m$ daje
$Dr_m+mhr_m=0$,
czyli $Dr_m/r_m=-mh$.
Wybór $u=r_m\in K^\times$ spełnia
żądaną tożsamość.
\end{proof}

\begin{przyklad}[Dołączenie wykładniczej]
\label{aa:prz:exponential-extension}
W $K=\mathbb C(x)$, $Dx=1$ nie istnieje
niezerowa funkcja wymierna $u$ spełniająca
$Du/u=c$ dla niezerowej stałej $c\in\mathbb C$.
Rzeczywiście, rozkład licznika i mianownika
na czynniki liniowe daje
\[
\frac{Du}{u}
=\sum_{a\in\mathbb C}\frac{m_a}{x-a},
\qquad m_a\in\mathbb Z,
\]
gdzie suma ma skończoną liczbę niezerowych wyrazów.
Po prawej jest funkcja właściwa, a więc
nie może być niezerową stałą.

Jeśli przestępny element $t$ spełnia
$Dt=t$, poprzedni lemat daje
$C_{K(t)}=\mathbb C$.
Analitycznie możemy wybrać $t=e^x$.
Przestępność tego zarodka można też wykazać
algebraicznie: gdyby $t$ był algebraiczny
nad $K$, skończony stopień $d=[K(t):K]$
i twierdzenie o śladzie i normie dałyby
\[
\frac{DN_{K(t)/K}(t)}{N_{K(t)/K}(t)}
=\operatorname{Tr}_{K(t)/K}(Dt/t)
=\operatorname{Tr}_{K(t)/K}(1)=d,
\]
co jest niemożliwe dla niezerowej stałej $d$
i funkcji wymiernej $N(t)\in\mathbb C(x)^\times$.
\end{przyklad}

\begin{definicja}[Rozszerzenie elementarne bez nowych stałych]
\label{aa:def:elementary-differential-extension}
Niech $(K,D)$ będzie ciałem różniczkowym
charakterystyki $0$.
\emph{Rozszerzeniem elementarnym} $E/K$
nazywamy rozszerzenie różniczkowe, dla
którego istnieje skończona wieża
$K=K_0\subseteq K_1\subseteq\cdots\subseteq K_r=E$,
gdzie $K_j=K_{j-1}(\theta_j)$, a każdy
generator $\theta_j$ spełnia co najmniej
jeden z trzech warunków:
\begin{enumerate}
\item $\theta_j$ jest algebraiczny nad $K_{j-1}$;
\item $D\theta_j=Du/u$ dla pewnego
      $u\in K_{j-1}^\times$ (krok logarytmiczny);
\item $\theta_j\ne0$ i
      $D\theta_j/\theta_j=Dv$ dla pewnego
      $v\in K_{j-1}$ (krok wykładniczy).
\end{enumerate}
Jeżeli ponadto $C_E=C_K$, mówimy, że
rozszerzenie jest \emph{bez nowych stałych}.
Dla $f\in K$ mówimy, że $f$ ma
\emph{elementarną funkcję pierwotną bez nowych stałych},
jeśli istnieje takie $E/K$ i $F\in E$
spełniające $DF=f$.
\end{definicja}

\begin{uwaga}[Granice definicji]
\label{aa:uw:elementary-versus-liouvillian}
Krok logarytmiczny ma postać
$D\theta=Du/u$, a \emph{nie} dowolną postać
$D\theta=a\in K_{j-1}$. Dopuszczenie dowolnej
funkcji pierwotnej już w definicji
prowadzi do większej klasy rozszerzeń
\emph{liouvillowskich}, rozważanych w
różniczkowej teorii Galois.
Wybór lokalnej gałęzi logarytmu lub
pierwiastka nie oznacza, że te funkcje
muszą istnieć jako jednowartościowe
funkcje globalne na całej płaszczyźnie.
Elementarną całkowalność będziemy
badać algebraicznie, a interpretować
lokalnie na wspólnym obszarze
określoności odpowiednich funkcji.
\end{uwaga}

\begin{przyklad}[Automorfizmy dwóch rozszerzeń]
\label{aa:prz:elementary-automorphisms}
Niech $K=\mathbb C(x)$ z $Dx=1$.

Dla $L=K(t)$, gdzie $t$ jest przestępny i $Dt=t$,
każdy automorfizm różniczkowy $\sigma:L\to L$
nad $K$ spełnia
$D(\sigma(t)/t)=0$, więc
$\sigma(t)=ct$ dla pewnego $c\in\mathbb C^\times$,
ponieważ $C_L=\mathbb C$.
Każdy taki wybór odwrotnie wyznacza automorfizm
o odwrotności $t\mapsto c^{-1}t$.
Grupa tych automorfizmów jest więc
izomorficzna z $(\mathbb C^\times,\cdot)$.

Dla $M=K(\ell)$, gdzie $\ell$ jest przestępny
i $D\ell=1/x$, mamy $C_M=\mathbb C$.
Równanie $D(\sigma(\ell)-\ell)=0$ daje
$\sigma(\ell)=\ell+c$ dla $c\in\mathbb C$,
a wszystkie te przesunięcia są automorfizmami.
Ich grupa jest izomorficzna z $(\mathbb C,+)$.
W obu przypadkach warunek komutowania z
derywacją ogranicza automorfizmy zwykłego
rozszerzenia ciał do bardzo konkretnej rodziny.
\end{przyklad}

Równości dla śladu i normy oraz kryteria
braku nowych stałych są dokładnie tymi
narzędziami, które pozwolą w następnym
rozdziale sprowadzić problem całkowania
elementarnego z wieży rozszerzeń do
tożsamości w początkowym ciele różniczkowym.
